\section{Preliminaries and the compressed formulation}
\label{sec:prelim}

\subsection{Basic notation}
\label{sec:prelim:notation}

Let $A\in\mathbb R^{n\times n}$, $C\in\mathbb R^{p\times n}$ and $L_0\in\mathbb R^{r_0\times n}$. The
design matrix is $\mathcal L=[L_0;R]\in\mathbb R^{q\times n}$, assumed to have full row rank, and its
row space is $V=\row\mathcal L\subseteq\mathbb R^n$. Put
\begin{equation}
  \label{eq:K0}
  K_0=\ker\left[\begin{smallmatrix}C\\ CA\end{smallmatrix}\right],\qquad
  k=\dim K_0,\qquad
  \rho=n-k,
\end{equation}
where $\rho=\rank[C;CA]$. Fix a basis $N\in\mathbb R^{n\times k}$ of $K_0$; Appendix~\ref{appB}
describes a normalized
choice. Here $\rho$ is the rank of the compressed data, whereas the layer index of this paper is $c$;
the two are different quantities. For a row space $V$ put
\begin{equation}
  \label{eq:compressed-data}
  H=LN,\qquad H_A=LAN,\qquad \omega=\rank H,\qquad c=k-\omega .
\end{equation}
The \emph{layer} is $c$. This paper treats only $c=1$, that is, $\omega=k-1$.

\subsection{The compressed criterion}
\label{sec:prelim:criterion}

The matrix $\mathcal L$ is feasible, that is, it satisfies \eqref{eq:5a} and \eqref{eq:5b}, if and only
if
\begin{equation}
  \label{eq:5ap}
  \row H_A\subseteq\row H,
\end{equation}
\begin{equation}
  \label{eq:5bp}
  \rank(\lambda H-H_A)=\omega\qquad\text{for all }\lambda\in\mathbb C .
\end{equation}
This equivalence, the identity $\omega=\rank[\mathcal L;CA;C]$, and the reformulation of
\eqref{eq:5bp} as the absence of a common root of the maximal minors of the pencil are established in
\cite{companion1}; they are used here and not re-proved. Condition \eqref{eq:5ap} only confines the
problem to a parameter space. What this paper adds to the criterion is a statement about \emph{which
points} a generic test covers, which is Lemma~\ref{sec4:lem:fullrank}.

\subsection{Standing assumptions}
\label{sec:prelim:assumptions}

Throughout the paper:
\begin{enumerate}[label=(\roman*),leftmargin=*,itemsep=2pt]
\item the standard Euclidean inner product is used, so restrictions to subspaces are automatically
nondegenerate;
\item input matrices are rational;
\item the design matrix $\mathcal L$ has full row rank, hence $\rank L_0=r_0$; otherwise no admissible
$\mathcal L$ exists at all, see Proposition~\ref{prop:rankprecheck};
\item the parameterization requires a basis of $K_0$; if $N'=NT$ with $T\in\operatorname{GL}_k(\mathbb R)$
is another basis, then $H'=HT$ and $H_A'=H_AT$, so $\lambda H'-H_A'=(\lambda H-H_A)T$. The change is
\emph{right} multiplication by an invertible matrix, which preserves ranks and hence feasibility.
\end{enumerate}

\begin{prop}
  \label{prop:rankprecheck}
  If $\rank L_0<r_0$, then for every $q$ there is no matrix $\mathcal L=[L_0;R]$ of full row rank, so
  the instance is infeasible.
\end{prop}

\begin{proof}
  Such a matrix would satisfy
  $q=\rank\mathcal L\le\rank L_0+(q-r_0)<r_0+(q-r_0)=q$, a contradiction.
\end{proof}

\subsection{Direction and the two functionals}
\label{sec:prelim:direction}

Let $V$ lie in the hyperplane layer. By \eqref{eq:compressed-data} and $c=1$ one has $\dim(V\cap K_0)=k-1$;
hence $\ker H$ is a line $\mathbb R v$ with $v\in K_0\setminus\{0\}$, and $V\subseteq U_v^\perp$,
where
\begin{equation}
  \label{eq:Uv}
  U_v=\operatorname{span}\{v,Av\}.
\end{equation}
Denote the two linear functionals on $K_0$ by
\begin{equation}
  \label{eq:functionals}
  \varphi_v(u)=\langle u,v\rangle,\qquad \varphi_{Av}(u)=\langle u,Av\rangle
  \qquad(u\in K_0).
\end{equation}
The inclusion $V\subseteq U_v^\perp$ is equivalent to saying that every $x\in V$ annihilates both $v$
and $Av$, that is, that $xN$ lies in the common kernel of the two functionals.

\subsection{Quantities used below}
\label{sec:prelim:quantities}

\begin{equation}
  \label{eq:compression}
  \begin{aligned}
    M&=\pi_{K_0}A|_{K_0},\qquad A_{21}=\pi_{K_0^\perp}A|_{K_0},\\
    E_\mu&=\ker(M-\mu I),\\
    E_\mu^0&=E_\mu\cap\ker(L_0N)\cap\ker(L_0AN),
  \end{aligned}
\end{equation}
\begin{equation}
  \label{eq:fibre-data}
  \begin{aligned}
    D_v&=K_0\cap U_v^\perp,\qquad \mathcal Q_v=U_v^\perp/D_v,\\
    V_0&=\row L_0,\qquad s=q-k+1,\\
    \delta&=\dim V_0-\dim(V_0\cap K_0),
  \end{aligned}
\end{equation}
together with the rank $r(v)=\dim\mathcal Q_v$ of the fibre. Its two values are computed in
Section~\ref{sec3}: $r(v)=\rho$ for eigen directions and $r(v)=\rho-1$ for non-eigen
directions.

\subsection{Parameter space and chart conventions}
\label{sec:prelim:charts}

The parameter space is taken to be a finite union of \emph{affine charts}: affine charts of projective
spaces for the direction, affine charts of Grassmannians for the fibre, and -- when
$\dim E_\mu^0>1$ -- affine charts of the direction base as well, combined with local trivializations of
the quotient bundle. For a real algebraic eigenvalue $\mu$ these charts are constructed over the number
field $K=\mathbb Q(\mu)$ equipped with the real embedding that selects $\mu$; their scalar extensions to
$\mathbb R$ are affine spaces whose real points are Zariski dense. All chart polynomials and
function-field gcd computations are performed over $K$. This convention is used in
Section~\ref{sec4} to pass from a generic algebraic point to a realizable real design.
